
\chapter{2021年上海卷（秋）}

\section{填空题}

\begin{problem}
已知 \(z_1=1+\i\)，\(z_2=2+3\i\)，求 \(z_1+z_2=\)\fillinblank{}.
\end{problem}

\begin{problem}
已知 \(A=\{x\mid 2x\le 1\}\)，\(B=\{-1,0,1\}\)，则 \(A\cap B=\)\fillinblank{}.
\end{problem}

\begin{problem}
若 \(x^2+y^2-2x-4y=0\)，求圆心坐标为\fillinblank{}.
\end{problem}

\begin{problem}
如图正方形 \(ABCD\) 的边长为 \(3\)，求 \(\overrightarrow{AB}\cdot \overrightarrow{AC}=\)\fillinblank{}.

\bitmapfigure[max width=6.0cm]{img/2021/shanghai/q4_fig1.png}
\end{problem}

\begin{problem}
已知 \(f(x)=\frac{3}{x}+2\)，则 \(f^{-1}(1)=\)\fillinblank{}.
\end{problem}

\begin{problem}
已知二项式 \((x+a)^5\) 展开式中，\(x^2\) 的系数为 \(80\)，则 \(a=\)\fillinblank{}.
\end{problem}

\begin{problem}
已知 \(x\le 3\)，\(2x-y-2\ge 0\)，\(3x+y-8\ge 0\)，\(z=x-y\)，则 \(z\) 的最大值为\fillinblank{}.
\end{problem}

\begin{problem}
已知 \(\{a_n\}\) 为无穷等比数列，\(a_1=3\)，\(a_n\) 的各项和为 \(9\)，\(b_n=a_{2n}\)，则数列 \(\{b_n\}\) 的各项和为\fillinblank{}.
\end{problem}

\begin{problem}
已知圆柱的底面圆半径为 \(1\)，高为 \(2\)，\(AB\) 为上底面圆的一条直径，C 是下底面圆周上的一个动点，则 \(\triangle ABC\) 的面积的取值范围为\fillinblank{}.
\end{problem}

\begin{problem}
已知花博会有四个不同的场馆 \(A\)，\(B\)，\(C\)，\(D\)，甲，乙两人每人选 \(2\) 个去参观，则他们的选择中，恰有一个馆相同的概率为\fillinblank{}.
\end{problem}

\begin{problem}
已知抛物线 \(y^2=2px\)（\(p>0\)），若第一象限的 \(A\)，\(B\) 在抛物线上，焦点为 \(F\)，\(|AF|=2\)，\(|BF|=4\)，\(|AB|=3\)，求直线 \(AB\) 的斜率为\fillinblank{}.
\end{problem}

\begin{problem}
已知 \(a_i\in \N^*\)（\(i=1\)，\(2\)，\(\dots\)，\(9\)），对任意的 \(k\in \N^*\)（\(2\le k\le 8\)），\(a_k=a_{k-1}+1\) 或 \(a_k=a_{k+1}-1\) 中有且仅有一个成立，\(a_1=6\)，\(a_9=9\)，则 \(a_1+\cdots+a_9\) 的最小值为\fillinblank{}.
\end{problem}

\section{单选题}

\begin{problem}
以下哪个函数既是奇函数，又是减函数（\quad）

\choices
    {\(y=-3x\)}
    {\(y=x^3\)}
    {\(y=\log_3 x\)}
    {\(y=3^x\)}
\end{problem}

\begin{problem}
已知参数方程 \(x=3t-4t^3\)，\(y=2t\sqrt{1-t^2}\)，\(t\in[-1,1]\)，以下哪个图符合该方程（\quad）

\choices
    {\choicebitmap[width=0.15\paperwidth]{img/2021/shanghai/q14_fig1.png}}
    {\choicebitmap[width=0.15\paperwidth]{img/2021/shanghai/q14_fig2.png}}
    {\choicebitmap[width=0.15\paperwidth]{img/2021/shanghai/q14_fig3.png}}
    {\choicebitmap[width=0.15\paperwidth]{img/2021/shanghai/q14_fig4.png}}
\end{problem}

\begin{problem}
已知 \(f(x)=3\sin x+2\)，对任意的 \(x_1\in\left[ 0,\frac{\pi}{2} \right]\)，都存在 \(x_2\in\left[ 0,\frac{\pi}{2} \right]\)，使得 \(f(x_1)=2f(x_2+\theta)+2\) 成立，则下列选项中，\(\theta\) 可能的值是（\quad）

\choices
    {\(\dfrac{3\pi}{5}\)}
    {\(\dfrac{4\pi}{5}\)}
    {\(\dfrac{6\pi}{5}\)}
    {\(\dfrac{7\pi}{5}\)}
\end{problem}

\begin{problem}
已知两两不相等的 \(x_1\)，\(y_1\)，\(x_2\)，\(y_2\)，\(x_3\)，\(y_3\)，同时满足 \circled{1}\(x_1<y_1\)，\(x_2<y_2\)，\(x_3<y_3\)；\circled{2}\(x_1+y_1=x_2+y_2=x_3+y_3\)；\circled{3}\(x_1y_1+x_3y_3=2x_2y_2\)，以下哪个选项恒成立（\quad）

\choices
    {\(2x_2<x_1+x_3\)}
    {\(2x_2>x_1+x_3\)}
    {\(x_2^2<x_1x_3\)}
    {\(x_2^2>x_1x_3\)}
\end{problem}

\section{解答题}

\begin{problem}
如图，在长方体 \(ABCD-A_1B_1C_1D_1\) 中，已知 \(AB=BC=2\)，\(AA_1=3\).

\begin{enumerate}
\item 若 \(P\) 是棱 \(A_1D_1\) 上的动点，求三棱锥 \(C-PAD\) 的体积；
\item 求直线 \(AB_1\) 与平面 \(ACC_1A_1\) 的夹角大小.
\end{enumerate}

\bitmapfigure[max width=5.2cm]{img/2021/shanghai/q17_fig1.png}
\end{problem}

\begin{problem}
在 \(\triangle ABC\) 中，已知 \(a=3\)，\(b=2c\).

\begin{enumerate}
\item 若 \(A=\frac{2\pi}{3}\)，求 \(S_{\triangle ABC}\)；
\item 若 \(2\sin B-\sin C=1\)，求 \(C_{\triangle ABC}\).
\end{enumerate}
\end{problem}

\begin{problem}
已知一企业今年第一季度的营业额为 \(1.1\) 亿元，往后每个季度增加 \(0.05\) 亿元，第一季度的利润为 \(0.16\) 亿元，往后每一季度比前一季度增长 \(4\%\).

\begin{enumerate}
\item 求今年起的前 \(20\) 个季度的总营业额；
\item 请问哪一季度的利润首次超过该季度营业额的 \(18\%\)？
\end{enumerate}
\end{problem}

\begin{problem}
已知椭圆 \(\Gamma:\frac{x^2}{2}+y^2=1\)，\(F_1\)，\(F_2\) 是其左，右焦点，直线 \(l\) 过点 \(P(m,0)\)（\(m\le -\sqrt2\)），交椭圆于 \(A\)，\(B\) 两点，且 \(A\)，\(B\) 在 \(x\) 轴上方，点 \(A\) 在线段 \(BP\) 上.

\begin{enumerate}
\item 若 \(B\) 是上顶点，\(\left| \overrightarrow{BF_1} \right|=\left| \overrightarrow{PF_1} \right|\)，求 \(m\) 的值；
\item 若 \(\overrightarrow{F_1A}\cdot \overrightarrow{F_2A}=\frac13\)，且原点 \(O\) 到直线 \(l\) 的距离为 \(\frac{4\sqrt{15}}{15}\)，求直线 \(l\) 的方程；
\item 证明：对于任意 \(m<-\sqrt2\)，使得 \(\overrightarrow{F_1A}\parallel \overrightarrow{F_2B}\) 的直线有且仅有一条.
\end{enumerate}
\end{problem}

\begin{problem}
已知 \(x_1\)，\(x_2\in \R\)，若对任意的 \(x_2-x_1\in S\)，\(f(x_2)-f(x_1)\in S\)，则有定义：\(f(x)\) 是在 \(S\) 关联的.

\begin{enumerate}
\item 判断和证明 \(f(x)=2x-1\) 是否在 \([0,+\infty)\) 关联？是否有 \([0,1]\) 关联？
\item 若 \(f(x)\) 是在 \(\{3\}\) 关联的，\(f(x)\) 在 \(x\in[0,3)\) 时，\(f(x)=x^2-2x\)，求解不等式 \(2\le f(x)\le 3\)；
\item 证明：\(f(x)\) 是 \(\{1\}\) 关联的，且是在 \([0,+\infty)\) 关联的，当且仅当“\(f(x)\) 在 \([1,2]\) 是关联的”.
\end{enumerate}
\end{problem}

